% formal/reeb-orbit-recovery.tex — orbit recovery material
% Sources: thesis/simple-minimizer-existence.tex

\begin{definition}[Simple Reeb orbit]\label{def:simple-reeb-orbit}
% Source: thesis/simple-minimizer-existence.tex
A generalized Reeb orbit $\gamma$ on a polytope $K$ is \emph{simple} if
its velocity is piecewise constant, each constant value is a facet Reeb
vector $R_i$, and for each facet index $i$ the set
\[
  \{ t \in [0,T] : \dot{\gamma}(t) = R_i \}
\]
is a single interval, possibly empty.
\end{definition}

\begin{lemma}[Base point recovery]\label{lem:base-point-recovery}
% Jörn: math approved (f7d26ea)
% Source: thesis/simple-minimizer-existence.tex
% Parameterization: K = {x : a_i^T x ≤ 1}, dual vertices a_i
Let $K = \{x : \langle a_i, x \rangle \leq 1\}$ be a polytope
and $(\sigma, \tau, b)$ a simple orbit
with $m$ active facets
$S = \{\sigma(k) : \tau_k > 0\}$.
Define the accumulated displacements
\[
  v_0 = 0, \qquad
  v_k = \sum_{j=1}^{k-1} \tau_j\, R_{\sigma(j)}
  \quad \text{for } k = 1, \ldots, m.
\]
A point $b \in \mathbb{R}^4$ is a valid base point
(i.e., the orbit $\gamma$ with $\gamma(0) = b$ and
velocity $R_{\sigma(k)}$ on the $k$-th segment
is a generalized Reeb orbit on~$\partial K$) if and only if:
\begin{enumerate}
\item \emph{On-facet (equality):}
  for each active $k$,
  $\langle a_{\sigma(k)},\, b \rangle
    = 1 - \langle a_{\sigma(k)},\, v_k \rangle$;
\item \emph{Inside $K$ (inequality):}
  for each active segment $k$ and all $s \in [0, \tau_k]$,
  $\langle a_j,\, b + v_k + s\, R_{\sigma(k)} \rangle
    \leq 1$ for all $j = 1, \ldots, F$.
\end{enumerate}
The equality constraints form
a linear system $A_S\, b = r$ with $m$ equations and $4$~unknowns.
The set of valid base points is a convex polytope of
dimension at most $4 - \operatorname{rank}(A_S)$.
\end{lemma}

\begin{proof}
\emph{Equality constraints.}
During the $k$-th segment, the orbit position is
\[
  \gamma(t) = b + v_k + (t - t_k)\, R_{\sigma(k)},
  \quad t \in [t_k,\, t_k + \tau_k],
\]
where $t_k = \sum_{j=1}^{k-1} \tau_j$.
For $\gamma$ to lie on facet $F_{\sigma(k)}$,
we require $\langle a_{\sigma(k)},\, \gamma(t) \rangle
= 1$ throughout the segment.
Since $R_{\sigma(k)}$ is tangent to $F_{\sigma(k)}$
($\langle a_{\sigma(k)}, R_{\sigma(k)} \rangle = 0$),
the time-dependent term vanishes:
\[
  \langle a_{\sigma(k)},\, \gamma(t) \rangle
  = \langle a_{\sigma(k)},\, b + v_k \rangle.
\]
Hence the on-facet condition for the entire segment
reduces to a single equation.

\emph{Inequality constraints.}
The orbit must remain inside~$K$:
$\langle a_j, \gamma(t) \rangle \leq 1$
for all $j$ and all~$t$.
Since $\gamma(t)$ is affine in~$t$ on each segment,
each inequality is linear in~$s$.
A linear function on $[0, \tau_k]$ attains its
maximum at an endpoint, so it suffices to
check at $s = 0$ and $s = \tau_k$ only.

\emph{Solution structure.}
The equality constraints are a linear system
$A_S\, b = r$ in~$b \in \mathbb{R}^4$.
Its solution set is an affine subspace of
dimension $4 - \operatorname{rank}(A_S)$.
The inequality constraints are linear in~$b$,
so the intersection is a convex polytope.
\end{proof}

\begin{remark}[Conversion from algorithm output]\label{rem:beta-to-tau}
% Jörn: math approved (f7d26ea)
% Source: thesis/simple-minimizer-existence.tex
% Parameterization: K = {x : a_i^T x ≤ 1}, dual vertices a_i
Algorithm~\ref{alg:ehz} outputs a permutation~$\sigma$ and
a vector~$\beta$ with $\beta_i \geq 0$ and
$\sum_i \beta_i = 1$.
The dwell times are
\[
  \tau_k = T\, \beta_{\sigma(k)},
\]
where $T = c_{\mathrm{EHZ}}(K)$ is the capacity
(the period of the orbit).
\end{remark}
